\documentclass[11pt]{article}
\usepackage[margin=1in]{geometry}
\usepackage{amsmath,amssymb}
\usepackage[hidelinks]{hyperref}
\setlength{\parindent}{0pt}
\setlength{\parskip}{0.5em}
\newcommand{\code}[1]{\texttt{\detokenize{#1}}}
\title{Student Audience Option: Mathematical Summary}
\author{}
\date{}
\begin{document}
\maketitle

\section{Modeling assumption}
The optional student mode represents an audience familiar with the building's
permitted evacuation paths. When enabled, every participant chooses a route
that minimizes remaining travel distance in the model's navigation graph.
This is an explicit behavioral assumption for the experiment. The option is
disabled by default, leaving the existing behavioral model in place.

Let the audience be indexed by $i\in\{1,\ldots,n\}$. Each student retains an
initial seat $s_i$, position $\mathbf x_i(t)$, floor, region, preferred speed,
reaction delay and body radius. Familiarity supplies knowledge of the building
layout; it does not change these physical characteristics.

\section{Distance representation and initial seat mapping}
Use the existing directed navigation graph
\[
 G=(V,E),\qquad L_e\geq0\quad(e\in E),
\]
whose vertices represent aisle ends, circulation connections, portal departures
and named destination entrances. Edge lengths are measured in metres.
Rows connect through their permitted aisle ends; stairs and access connectors
use their declared traversal lengths. Circulation distances follow the
schematic rectangular or curved region geometry.

Every seat receives an initial shortest route to a permitted exterior
threshold. The graph preserves floor and region connectivity, so a short
straight line through a wall, seating block or floor is not an admissible
route. The elevator remains excluded. These are distances in the implemented
schematic graph, rather than surveyed three-dimensional building distances.

\section{Student route-choice rule}
Let $B_i(t)$ be the set of closures known to student $i$. Remove the associated
transfer edges and exterior-threshold destinations from the route search.
For a graph state $v$, define
\[
 D_{B_i(t)}(v)=
 \min_{\pi:v\leadsto V_{\mathrm{exit}}(B_i(t))}
                         \sum_{e\in\pi}L_e,
\]
where $V_{\mathrm{exit}}(B_i(t))$ contains the exterior thresholds not known
to be closed. Set this distance to infinity if no admissible route exists.
The implementation evaluates these distances using reverse multi-source
Dijkstra search and caches fields for repeated closure-knowledge states.

At a decision time, let $\mathcal A_i^{\mathrm{phys}}(t)$ contain the student's
physically permitted next-target actions. For action $a$, let $L_i(a,t)$ be
the distance to and through that action, and $v_a$ its downstream graph state.
A portal action includes its actual transfer before its continuation is scored.
The remaining distance associated with the action is
\[
 Q_i(a,t)=L_i(a,t)+D_{B_i(t)}(v_a).
\]
The student selects
\[
 \boxed{\displaystyle
 a_i^*(t)\in\arg\min_{a\in\mathcal A_i^{\mathrm{phys}}(t)}Q_i(a,t).}
\]
Unreachable actions are excluded. Equal distances use deterministic exit-ID
and action-ID tie rules. If every continuation is unreachable, the student
clears its current target and waits for reconsideration.

The initial seat mapping provides a reference route. Subsequent decisions
use the current position and knowledge, allowing the intended route to change
after physical displacement or newly discovered closures. The existing default
three-second decision interval is retained, with immediate reconsideration
on relevant region transitions and local closure discovery.

\section{Knowledge and unexpected obstructions}
Students know target locations and connectivity, so familiar targets can be
considered without the ordinary target-visibility filter. Physical eligibility
still restricts available actions to the correct floor, region and connectors.

Layout knowledge does not imply immediate knowledge of an unexpected blocked
exit. The distance query uses $B_i(t)$ rather than all physically disabled
portals. Students retain the existing visibility-based local discovery rule,
with a default discovery radius of 2.2 metres. A student may therefore approach
a closed door, discover its condition, and then choose the shortest remaining
known route. Physical closure prevents passage even before discovery.

The directed movement model defines which alternatives exist. If the current
state has no permitted alternate continuation after a closure, the policy waits
rather than inventing a reverse movement or an inter-floor shortcut.

\section{Relationship to behavioral utility and physical movement}
Student route selection minimizes distance directly. Behavioral weights,
peer-following preference, norm penalties, target-switching penalties,
perceived queue time and panic/residual utility shocks do not select the route.
The existing four weights are still initialized for consistent audience
construction, but are unused by the student route-choice rule. This option
does not add a familiarity weight to the ordinary behavioral utility.

After selecting a target, the student moves through the existing physical
engine. In manuscript force mode the underlying acceleration includes
\[
 \frac{d\mathbf v_i}{dt}
 =\frac{v_i^0\widehat{\mathbf e}_i-\mathbf v_i}{\tau}
  +\mathbf f_i^{\mathrm{obs}}
  +\sum_j\mathbf f_{ij}^{\mathrm{int}},
\]
where $\widehat{\mathbf e}_i$ points toward the legal waypoint for the selected
action. Configured integration, local avoidance, collision exclusion, region
constraints, door capacities and stair queues remain active. Consequently,
shortest intended distance does not guarantee the shortest realized trajectory
or the fastest evacuation through a congested building.

\section{Precedence over crowd managers}
When student mode is active,
\[
 \mathcal M_{\mathrm{active}}=\varnothing.
\]
No crowd-manager bodies, instructions, authority bonuses or notifications
affect the audience, even if the manager option is selected. Explicit physical
closure schedules remain applicable; staff-notification delays have no effect
because there are no staff. The GUI disables the manager control while
retaining its preference for use after student mode is switched off and the
experiment is reset.
The optional LED-strip mode is also suppressed: students receive no LED
instruction or utility bonus and no strip network is displayed. Its controls
retain their preferences for a later nonstudent reset. Student route choice
continues to use its own direct distance-minimization rule.

\section{Controls and implementation}
Enable \emph{Students: shortest routes (apply/reset)} in the GUI and select
\emph{Reset / apply controls}. The command-line equivalent is
\begin{verbatim}
python run.py --n 300 --students --gui
\end{verbatim}
World JSON can instead contain \code{"studentMode": true}.
The \code{--no-students} flag restores ordinary behavioral route choice.

\code{students.py} implements the student policy and reuses the graph in
\code{crowd_routes.py}. Saved student decisions identify their selection mode,
chosen action, intended exit and remaining route distances. Unreachable values
are represented as nulls in JSON. These diagnostics make the assumed
distance-minimizing behavior inspectable without changing physical movement.
\end{document}
